\documentclass[10pt,a4paper]{amsart}
\usepackage[margin=19mm]{geometry}
\usepackage{amsmath,amssymb,mathtools}
\usepackage[scr=rsfs,cal=euler]{mathalfa}
\usepackage{xfrac}
\usepackage[unicode,hidelinks,bookmarks=false]{hyperref}
\newcommand{\ie}{i.e.\ }
\newcommand{\cf}{cf.\ }
\newcommand{\resp}{respectively\ }
\newcommand{\Subsection}{\subsection}
\providecommand{\nobreakdash}{\nobreak\hbox{-}}
\setlength{\emergencystretch}{2em}
\multlinegap0pt
\usepackage{bm}
\usepackage{bbm}
\usepackage{dsfont}
\usepackage{xspace}
\usepackage{cancel}
\usepackage{mathtools}
\usepackage{cleveref}
\usepackage{amsthm}
\makeatletter
\@ifundefined{lemma}{\newtheorem{lemma}{Lemma}}{}
\@ifundefined{theorem}{\newtheorem{theorem}{Theorem}}{}
\makeatother
\makeatletter
\def\FF{\mathbb{F}}
\def\RR{\mathbb{R}}
\def\ZZ{\mathbb{Z}}
\expandafter\ifx\csname breakablealgorithm\endcsname\relax
\newenvironment{breakablealgorithm}
  {% \begin{breakablealgorithm}
   \begin{center}
     \refstepcounter{algorithm}% New algorithm
     \hrule height.8pt depth0pt \kern2pt% \@fs@pre for \@fs@ruled
     \renewcommand{\caption}[2][\relax]{% Make a new \caption
       {\raggedright\textbf{\ALG@name~\thealgorithm} ##2\par}%
       \ifx\relax##1\relax % #1 is \relax
         \addcontentsline{loa}{algorithm}{\protect\numberline{\thealgorithm}##2}%
       \else % #1 is not \relax
         \addcontentsline{loa}{algorithm}{\protect\numberline{\thealgorithm}##1}%
       \fi
       \kern2pt\hrule\kern2pt
     }
  }{% \end{breakablealgorithm}
     \kern2pt\hrule\relax% \@fs@post for \@fs@ruled
   \end{center}
  }
\fi
\newcommand{\setdes}[2]{\left\{#1\;\middle|\;#2\right\}}
\newcommand{\sgn}{\mathop{\rm sgn}\nolimits}
\makeatother
\title[Engineered Complete Intersections: Algorithmic Aspects]{Engineered Complete Intersections: Algorithmic Aspects}
\author{Alexander Esterov, Rafael Mohr, Yulia Mukhina}
\date{Source: arXiv:2607.23622v1 (CC BY 4.0)}
\begin{document}
\maketitle

\noindent\emph{Source and licence.} Engineered Complete Intersections: Algorithmic Aspects. Version arXiv:2607.23622v1 (CC BY 4.0), distributed under the Creative Commons Attribution 4.0 International licence, as recorded in the arXiv licence metadata for this exact version and shown on the abstract page https://arxiv.org/abs/2607.23622v1. Full rights notice: licenses/arxiv-2607.23622-CC-BY-4.0.txt.\par\smallskip

\noindent\textit{Editorial scope.} Counted unit: the Theorem whose source label is \texttt{thm:rootcounting} in the article (source0.tex lines 1357--1371, source label \texttt{thm:rootcounting}), delivered together with the article's own complete original proof of it (source0.tex lines 1372--1493, one \texttt{proof} environment of 122 lines). No mathematics is written, repaired or re-derived: statement and proof are the article's own text line for line. Required context retained uncounted: eq:polsys (lines 1251--1254); ltriv (lines 1288--1291). Every definition, display and label of those lines is retained unchanged. Omitted from this package and not redistributed: 1--1250, 1255--1287, 1292--1356, 1494--2646. Nothing inside a retained excerpt is omitted or altered: every retained body is delivered exactly as the article's lines are written, so no omission receipt of any kind accompanies this package. Citations: the retained excerpts carry no citation command at all, so no bibliography entry of the article is transcribed and no reference is redistributed. Reference adaptation: none was needed and none was made; every reference inside the delivered unit resolves inside it, so no unresolved reference or citation is produced. Printed slips kept literal: no printed word, symbol, hypothesis or number of the counted statement or of its proof was altered. Layout and class: the article's own class is \texttt{article} with packages including geometry, authblk, bm, float, fontenc, changepage, caption, algpseudocodex, parskip, booktabs, xcolor, algorithm, graphicx, enumitem; the wrapper is the lane's pinned \texttt{amsart} editorial wrapper at 10pt on A4, which restates the theorem-like environments the retained text needs and adds the article's own macro definitions that the retained text needs (\texttt{\textbackslash FF}, \texttt{\textbackslash RR}, \texttt{\textbackslash ZZ}, \texttt{\textbackslash setdes}, \texttt{\textbackslash sgn}), with extra wrapper packages (bm, bbm, dsfont, xspace, cancel, mathtools, cleveref). The delivered numbering is the unit's own. Assets: no retained span contains a figure environment, a graphics-inclusion command, a PDF-page inclusion, a table or a TikZ picture; no image file is copied into this package, the package declares no asset dependency and no image file is redistributed with it. Licence: version arXiv:2607.23622v1 of this article is distributed under the Creative Commons Attribution 4.0 International licence, as recorded in the arXiv licence metadata for this exact version and shown on the abstract page https://arxiv.org/abs/2607.23622v1. A full rights notice is delivered at licenses/arxiv-2607.23622-CC-BY-4.0.txt. Characters: every retained excerpt is pure ASCII, so the delivered engine, XeLaTeX with its default Latin Modern text font, reports no missing character.
\begin{equation}
  \label{eq:polsys}
  f_{ij}(x)=0, \qquad i=1,\ldots,k,\quad j=1,\ldots,d_i,
\end{equation}

\begin{lemma}\label{ltriv}
  Once any maximal minor of any of the matrices $U_i$ is zero, the
  system of equations \eqref{eq:polsys} has no solutions in $(\RR^{*})^n$.
\end{lemma}

\begin{theorem}\label{thm:rootcounting}
  Assume that the system of equations \eqref{eq:polsys} does not
  satisfy the conditions of \Cref{ltriv}. Under the identification of
  $\FF_2^n$ with the group of orthants $E$ given above, we have
  \[\setdes{e\in E}{\mathcal{O}_e\text{ contains a solution of \eqref{eq:polsys}}} \cong \Lambda^{-1}(b).\]

  Consequently, the total number of orthants with a solution of
  \eqref{eq:polsys} (and thus the total number of real toric solutions
  for square systems of the shape \eqref{eq:polsys}) not satisfying the
  condition of \Cref{ltriv}) equals
  \begin{itemize}
  \item $0$ if $b$ is not contained in the image of $\Lambda$;
  \item $2^{n-\operatorname{rank}_{\mathbb F_2}\Lambda}$ otherwise.
  \end{itemize}
\end{theorem}

\begin{proof}
  By assumption, every coordinate of every cofactor vector $c_i$ is
  nonzero. Fix $e\in E$ and define
  \begin{align*}
    q_\ell&:=(-1)^{e_\ell}\\
    q&:=(q_1,\ldots,q_n),
  \end{align*}
  and consider the corresponding orthant $\mathcal O_e$. Every
  $x\in\mathcal O_e$ can be written uniquely as $x = q\cdot y$ for some
  $y\in\RR_{>0}^n$ under coordinatewise multiplication. Then, for every
  $a\in \ZZ^n$, we have
  \[x^a = q^ay^a\quad \text{ where }\quad q^a=(-1)^{\sum_{\ell=1}^n a_\ell e_\ell}.\]

  For every block of \eqref{eq:polsys} indexed by $S_i$, the equations
  $f_{ij}(x)=0$, $j=1,\ldots,d_i$, are equivalent to
  $$
  \left(\sigma_i(a)x^a\right)_{a\in S_i}\in\ker U_i=\RR c_i.
  $$
  Hence if $x$ is a solution of \eqref{eq:polsys} then for every $i$
  there is a sign $\eta_i\in\{\pm1\}$ such that
  $$
  \sigma_i(a)q^a
  =
  \eta_i\sgn c_i(a)
  $$
  for every $a\in S_i$. Equivalently,
  $$
  q^a\sgn\bigl(\sigma_i(a)c_i(a)\bigr)
  $$
  is independent of $a\in A_i$.
  By the definition of $b_i$, this says that the function
  $$
  a\longmapsto \left[\sum_{\ell=1}^n a_\ell e_\ell+b_i(a)\right]
  $$
  is constant on $A_i$. Since addition and subtraction coincide in
  $\mathbb F_2$, this is equivalent to
  $$
  \left[
    a\longmapsto\sum_{\ell=1}^n a_\ell e_\ell
  \right]
  =
  [b_i].
  $$
  Requiring this for every $i$ gives
  $$
  \Lambda(e)=b.
  $$
  Thus every solution-containing orthant of $\RR^n$ yields an
  element of $\Lambda^{-1}(b)$.

  Conversely, suppose that $\Lambda(e)=b$. Then, for every $i$, the function
  $$
  a\longmapsto \left[\sum_{\ell=1}^n a_\ell e_\ell+b_i(a)\right]
  $$
  is constant on $S_i$. Consequently, there is a sign
  $\eta_i\in\{\pm1\}$ such that
  $$
  s_i(a)q^a
  =
  \eta_i\sgn c_i(a)
  $$
  for every $a\in S_i$. It remains to construct the absolute value of
  the root corresponding to $e$. Consider the linear map
  \begin{align*}
    &L:\RR^n\longrightarrow\RR^D,\\
    &L(v) := \left(
    \langle a_{ik}-a_{i0},v\rangle\right)_{\substack{1\leq i\leq r\\1\leq k\leq d_i}}.  
  \end{align*}
  The vectors $a_{ik}-a_{i0}$ are pairwise linearly independent by the
  transversality assumption on the $S_i$. Therefore $L$ has rank $D$ and
  is surjective. We can thus choose $v\in\RR^n$ satisfying
  $$
  \langle a_{ik}-a_{i0},v\rangle
  =
  \log\frac{|c_i(a_{ik})|}{|c_i(a_{i0})|}
  $$
  for every $i$ and every $k=1,\ldots,d_i$. Put
  $$
  y_\ell=e^{v_\ell},
  \qquad
  \ell=1,\ldots,n.
  $$
  Then, for every $i$, there is a positive number $\mu_i$ such that
  $$
  y^a=\mu_i|c_i(a)|
  $$
  for all $a\in S_i$. Finally setting $x:=q\cdot y\in\mathcal O_e$, we obtain
  $$
  \sigma_i(a)x^a
  =
  \sigma_i(a)q^ay^a
  =
  \eta_i\mu_i c_i(a).
  $$
  Hence
  $$
  \left(\sigma_i(a)x^a\right)_{a\in S_i}
  =
  \eta_i\mu_i c_i\in\ker U_i,
  $$
  and so $x$ is a solution of \eqref{eq:polsys} contained in $\mathcal O_e$. We have
  thus finally proven
  $$
  e\in\Lambda^{-1}(b)
  \quad\Longleftrightarrow\quad
  \mathcal O_e\text{ contains a solution of \eqref{eq:polsys}}.
  $$
  Because $e\mapsto\mathcal O_e$ is a bijection from $\FF_2^n$ to the group
  of orthants of $\RR^n$, it restricts to a bijection from the fiber
  $\Lambda^{-1}(b)$ to the set of orthants containing a solution of
  \eqref{eq:polsys}.

  If $b\notin\operatorname{im}\Lambda$, this fiber is empty. Otherwise it is an
  affine space in $\FF_2^n$ of dimension
  $\dim\ker \Lambda = n-\operatorname{rank} \Lambda$. Such an affine
  space has cardinality $2^{n-\operatorname{rank}_{\mathbb F_2}\Lambda}$. 

  Finally, if $D=n$, the map $L$ is an isomorphism. Thus the vector
  $v$, and hence the solution $x$ constructed above in a prescribed orthant, is
  unique. Therefore, for square systems, the number of real toric solutions equals
  the number of solution-containing orthants.
\end{proof}

\end{document}
